%! Author = wboughattas
%! Date = 2025-09-01

% Preamble
\documentclass[11pt]{article}

% Packages
\usepackage{amsmath}

% Document
\begin{document}
    \section{Recap: Finance and Machine Learning Primer}

    This section provides a self-contained guide to the finance and ML terminology used in the plan.
    It is written in plain English, with definitions and small numerical examples.
    Readers may skim headings they know and focus on unfamiliar ones.

    \subsection{The Absolute Basics}

    \paragraph{Candles, Bars, OHLCV.}
    A \emph{bar} or \emph{candle} is one time bucket of market data (e.g.\ 1 minute, 1 day).
    OHLCV stands for \emph{Open, High, Low, Close, Volume}.
    \begin{itemize}
        \item Example: a 1-minute bar may have Open=100.00, High=100.20, Low=99.95, Close=100.10, Volume=50{,}000.
    \end{itemize}

    \paragraph{Horizon $n$.}
    The forecast horizon. If $n=5$ bars on 1-minute data, we are predicting 5 minutes ahead.

    \paragraph{Returns (simple vs.\ log).}
    \begin{align*}
        R &= \frac{P_{t+n} - P_t}{P_t} \quad \text{(simple return)} \\
        r &= \ln\!\left(\frac{P_{t+n}}{P_t}\right) \quad \text{(log return)}
    \end{align*}
    For small moves, $r \approx R$. Log returns are preferred because they add over time and work well in statistics.
    Example: $P_t=100$, $P_{t+n}=101 \Rightarrow R=1\%$.

    \paragraph{Basis Points (bps).}
    1 bp = 0.01\% = 0.0001 in decimal.
    Examples: 5 bps = 0.05\%, 25 bps = 0.25\%.

    \subsection{Uncertainty, Risk, and ``Edge''}

    \paragraph{Expected Return $\mu$ and Uncertainty $\sigma$.}
    \begin{itemize}
        \item $\mu_n(t)$: model’s expected log return over next $n$ bars.
        \item $\sigma_n(t)$: predicted uncertainty (often a standard deviation).
        \item Example: for next 5 bars, $\mu=+15$ bps, $\sigma=20$ bps.
    \end{itemize}

    \paragraph{Edge.}
    Expected gain after accounting for costs and noise.
    \[
        U_n(t) = |\mu_n(t)| - c(t) - \kappa \cdot \sigma_n(t),
    \]
    where $c(t)$ = trading cost, $\kappa$ = risk buffer coefficient.
    Absolute value ensures both positive (long) and negative (short) predictions can be attractive if large enough.

    \paragraph{Intuition.}
    If the predicted move barely beats costs, skip it. Trade only when the edge is comfortably above costs and noise.

    \subsection{Trading Costs and Execution}

    \paragraph{Commission.}
    Broker fee per trade, often quoted in bps or cents/share.

    \paragraph{Spread and Slippage.}
    \begin{itemize}
        \item Bid--ask spread: buyers bid 99.99, sellers ask 100.01. Spread = 2 cents.
        \item Crossing the spread: market order pays $\tfrac{1}{2}$ spread to enter and $\tfrac{1}{2}$ to exit.
        \item Slippage: any adverse price move between decision and fill.
    \end{itemize}
    Rule of thumb: assume at least half-spread cost per market trade.

    \paragraph{Shorting.}
    To short, borrow shares and sell them. May incur a borrow fee. Some stocks are not shortable.

    \paragraph{Latency and Fill Timing.}
    Backtests must specify: do we execute at next bar’s open or close? Consistency is crucial.

    \subsection{Position Sizing, Exits, and Risk Controls}

    \paragraph{Position Sizing.}
    Scale by conviction relative to risk:
    \[
        S_t = \text{clip}\!\left(\beta \cdot \frac{\mu}{\sigma}, -S_{\max}, S_{\max}\right).
    \]
    \begin{itemize}
        \item $\beta$: aggressiveness (typical range 0.3--1.0).
        \item $S_{\max}$: cap on exposure (e.g.\ $\pm 1.0$).
        \item If $\mu/\sigma$ small, take small size; if large, take larger (capped).
    \end{itemize}

    \paragraph{Exit Rules.}
    \begin{itemize}
        \item \emph{Time stop:} exit after $n$ bars if no exit occurred.
        \item \emph{Stop-loss (SL):} exit if loss exceeds threshold (e.g.\ $1\sigma$).
        \item \emph{Take-profit (TP):} exit if profit exceeds threshold.
    \end{itemize}

    \paragraph{Turnover.}
    How often positions change. High turnover $\Rightarrow$ higher costs.

    \paragraph{Drawdown (DD).}
    Peak-to-trough equity loss.
    Example: account rises to \$1{,}000, falls to \$900, then recovers. Max DD = \$100 = 10\%.

    \paragraph{Exposure.}
    \begin{itemize}
        \item Net exposure = long $-$ short. Example: +60\% long, $-$20\% short $\Rightarrow$ net +40\%.
        \item Gross exposure = sum of longs and shorts. Example: 60\%+20\%=80\%.
    \end{itemize}

    \paragraph{Capacity.}
    The maximum capital a strategy can handle before costs erode profits. Test by scaling trades and observing net P\&L.

    \subsection{Model Predictions and Calibration}

    \paragraph{Probability vs.\ Regression.}
    \begin{itemize}
        \item Classification: predict $p_{\text{up}}$ (probability of up move).
        \item Regression: predict $\mu$ (expected return), sometimes $\sigma$ (uncertainty).
    \end{itemize}

    \paragraph{Calibration.}
    Raw probabilities may be overconfident.
    \begin{itemize}
        \item Platt scaling: logistic mapping.
        \item Isotonic regression: flexible monotonic mapping.
    \end{itemize}

    \paragraph{Quantile Regression.}
    Predict quantiles (e.g.\ 20th, 50th, 80th percentiles). Trained via pinball loss. Gives prediction intervals and supports position sizing.

    \subsection{Evaluation Metrics}

    \paragraph{Sharpe Ratio.}
    \[
        \text{Sharpe} = \frac{\text{mean excess return}}{\text{standard deviation of returns}}.
    \]
    Annualized by $\sqrt{252}$ for daily data. Measures reward per unit risk.

    \paragraph{Sortino Ratio.}
    Like Sharpe, but only penalizes downside volatility.

    \paragraph{Calmar Ratio.}
    Annual return divided by max drawdown.

    \paragraph{Hit Rate and Average Trade P\&L.}
    Hit rate = fraction of profitable trades.
    Average trade P\&L = mean profit per trade after costs.

    \paragraph{Calibration Metrics.}
    \begin{itemize}
        \item Brier score: squared error between predicted probabilities and outcomes.
        \item Reliability curve: compares predicted probability bins with actual outcomes.
    \end{itemize}

    \subsection{Data Hygiene and Validation}

    \paragraph{Look-ahead Leakage.}
    Using data not yet available at decision time (e.g.\ earnings before release). Prevent via timestamping and lagging.

    \paragraph{Corporate Actions.}
    Splits, dividends, and symbol changes must be adjusted for consistent returns.

    \paragraph{Purged, Embargoed K-Fold.}
    For time-series, classic K-fold leaks. Purging removes overlapping samples; embargo adds buffer gaps.

    \paragraph{Walk-Forward Testing.}
    Train on window A, validate on B, test on C. Roll forward. Mimics live deployment.

    \paragraph{Multiple Testing.}
    Trying many models risks overfitting. Keep a final untouched test.

    \subsection{Labeling and Meta-Labeling}

    \paragraph{Labels.}
    Direction: $y_t=1$ if $P_{t+n}>P_t$, else 0. Regression: target $r_{t\to t+n}$.

    \paragraph{Triple-Barrier Method.}
    Define three exits: take-profit, stop-loss, time. Label is which barrier hits first.

    \paragraph{Meta-Labeling.}
    Base model proposes trades, meta-model predicts whether to act. Often reduces churn and improves Sharpe.

    \subsection{Optimization Methods}

    \paragraph{Bayesian Optimization.}
    Efficient hyper-parameter tuning via surrogate modeling.

    \paragraph{End-to-End Training.}
    Optimize directly through backtest metrics. Hard due to non-differentiability. Approaches: black-box threshold tuning, short episodes, or soft approximations.

    \subsection{Microstructure and Regimes}

    \paragraph{Liquidity.}
    Ease of trading without moving price. High liquidity $\Rightarrow$ tight spreads.

    \paragraph{Regimes.}
    Market ``moods'' such as high-vol vs low-vol, trending vs mean-reverting.

    \paragraph{Volatility of Volatility.}
    How fast volatility itself changes. High vol-of-vol $\Rightarrow$ noisy signals.

    \subsection{Artifacts and Outputs}

    \paragraph{Equity Curve.}
    Cumulative P\&L over time (start at 1.00, plot growth).

    \paragraph{Blotter.}
    Trade log: timestamp, side, size, price, costs, entry/exit reasons.

    \paragraph{Gross vs Net P\&L.}
    Gross = before costs. Net = after costs. A strategy with good gross but poor net is cost-dominated.

    \subsection{Worked Examples}

    \paragraph{Example 1 (Skip Due to Costs).}
    $\mu=+7$ bps, $\sigma=15$ bps, $c=4$ bps, $\kappa=0.5 \Rightarrow$ buffer=7.5 bps. Net edge=3 bps $<$ 7.5. Skip trade. Price rises 5 bps, but costs would erase gains.

    \paragraph{Example 2 (Take 5-Bar Long).}
    $\mu=+18$ bps, $\sigma=20$ bps, $c=3$ bps, $\kappa=0.5$. Net edge=15 bps $>$ 10 bps buffer. Position size = $0.8 \times (18/20)=0.72$. Price rises 22 bps. P\&L positive.

    \paragraph{Example 3 (Short on Longer Horizon).}
    1-bar: $\mu=+3$, $\sigma=12$, $c=3 \Rightarrow U=-6$ bps (skip).
    20-bar: $\mu=-40$, $\sigma=25$, $c=3 \Rightarrow U=24.5$ bps (trade).
    Take short of $-1.0$. Price drifts down; profit realized.

    \paragraph{Example 4 (Meta-Label Blocks Trade).}
    Edge barely clears buffer, but spreads widening. Meta-model says low tradeability. Skip. Market chops; skipping avoids costs.

    \paragraph{Example 5 (Controlled Loss).}
    Strong short signal. Two bars later, market rallies +30 bps. Stop-loss triggers. Small capped loss avoids larger drawdown.

    \subsection{Prerequisites Checklist}

    \begin{itemize}
        \item Basic algebra and logs (why log returns add).
        \item Mean and standard deviation as measures of uncertainty.
        \item Bid/ask mechanics and spread as cost.
        \item Time-series validation (purge/embargo, walk-forward).
        \item Probability calibration.
    \end{itemize}


\end{document}